\chapter{Background \& Literature Review}

Re-LAION-5B comprises approximately 5.85 billion image-text pairs harvested via Common
Crawl~\cite{Schuhmann_2022_Laion5b}, re-released with improved safety
filtering~\cite{Laion_2024_Relaion5B}; its LAION-Aesthetics v2~5+ subset constitutes
the training corpus for SDv1.5~\cite{Schuhmann_2022_LaionAesthetics}. SDv1.5 is a
Latent Diffusion Model~\cite{Rombach_2022_CVPR} that denoises in a compressed latent
space via a DDPM U-Net conditioned on CLIP text
embeddings~\cite{Ho_2020_DDPM,Radford_2021_CLIP}; higher classifier-free guidance scales
amplify demographic stereotypes~\cite{Ho_2022_CFG}. Its fully documented training lineage
makes it uniquely suited for mechanistic bias auditing, unlike later models such as
SDXL~\cite{Podell_2023_SDXL} that rely on proprietary corpora. Zeng et
al.~\cite{Zeng_2024_UnderstandingBiasinLargeScaleVisualDatasets} operationalise spurious
feature analysis through a transformation-based classification framework in which each
transformation, spanning pixel shuffle, patch shuffle, mean RGB, frequency filtering,
Canny edge, SAM contour, monocular depth, semantic segmentation, object detection, and
image captioning, isolates a distinct visual property while discarding all others,
revealing that dataset identity is redundantly encoded across colour, texture, geometry,
and semantic content simultaneously and that generative models reproduce training data
biases in their outputs. Meister et
al.~\cite{Meister_2023_GenderArtifactsinVisualDatasets} extend this to demographic
attributes via a systematic occlusion-based probing protocol, establishing that gender
artifacts are distributed across the entire image structure rather than localised to the
subject, undermining fairness-through-blindness assumptions. This thesis extends Zeng et
al.'s framework to a three-class Re-LAION-5B / SDv1.5 / COCO setting and applies
Meister et al.'s occlusion protocol to both retrieved and generated corpora.

Occupation prompts drawn from the ISCO-08 taxonomy~\cite{ILO_ISCO08} provide an
interpretable probe for demographic
priors~\cite{Naik_2023_SocialBiasesThroughTheTextToImageGenerationLens,
Bianchi_2023_TTIGenAmplifiesDemographicStereotypes,
Luccioni_2023_AnalyzingSocietalRepresentationsDiffusionModels}. Demographic annotation
uses automated classifiers applied to segmented facial skin regions to prevent background
shortcut learning~\cite{Matias_2024_EnhancingFairnessMST}, with skin tone operationalised
via the ten-point Monk Skin Tone scale~\cite{Schumann_2024_MonkSkinTone} and evaluated
against CCv2~\cite{Porgali_2023_CCv2} and FACET~\cite{Gustafson_2023_FACET}; automated
annotation introduces structured measurement error~\cite{Buolamwini_2018_GenderShades}.
Inference-time debiasing is applied via ITI-GEN~\cite{Zhang_2023_ITIGEN}, which learns
inclusive text-space tokens grounded in demographic reference images, composed with
ControlNet OpenPose~\cite{LZhang_2023_ControlNet} for structural conditioning and
Chamfer-based colour alignment~\cite{Shum_2025_ColorAlignment} to neutralise colour-level
spurious priors.
